\documentclass[11pt,a4paper]{article}
\usepackage[margin=2.4cm]{geometry}
\usepackage{amsmath,amssymb}
\usepackage[T1]{fontenc}
\usepackage{lmodern}
\usepackage{microtype}
\usepackage{enumitem}
\usepackage{hyperref}
\setlist{itemsep=2pt,topsep=3pt}

\title{Report on \emph{``Matrix Theory Reloaded: A BPS Road to Holography''}\\[2pt]
\large by C.~D.~A.~Blair, J.~Lahnsteiner, N.~A.~Obers and Z.~Yan}
\date{}

\begin{document}
\maketitle
\vspace{-1.5em}

\paragraph{Summary of the paper.}
The paper treats the decoupling limit that produces matrix theory on D$p$-branes as a
\emph{BPS decoupling limit} of the full type~II theory. In this limit the target space
becomes a non-Lorentzian ``matrix $p$-brane theory'' (M$p$T) geometry with a
codimension-$(p+1)$ foliation and a critical RR potential. The paper makes three main claims.
(i)~Applying the M$p$T limit at asymptotic infinity of the D$p$-brane supergravity
solution produces exactly the Maldacena/IMSY near-horizon limit. Applying a second, misaligned
limit (M$p$T on a D$q$-brane background) systematically produces non-Lorentzian bulk geometries.
These reproduce the recent proposals of Lambert et al.\ and of Fontanella et al.
(ii)~After uplift to M-theory, each BPS limit is a DLCQ. This motivates a conjecture that
holography is a DLCQ$^n$/DLCQ$^m$ correspondence with $m>n$: AdS$_5$/CFT$_4$ is
DLCQ$^0$/DLCQ$^1$, and D1--D5 AdS$_3$/CFT$_2$ is DLCQ$^0$/DLCQ$^2$.
(iii)~Undoing the limit is a $p$-brane generalisation of $T\bar T$. Promoting $\omega\to\mathbb{\Omega}(r)$
turns the flat asymptotic M$p$T geometry into the near-horizon bulk. On the worldvolume this gives
closed-form flow equations that take $D$ free scalars in $p+1$ dimensions to the
Dirac--Nambu--Goto action for every $p$, plus exact Maxwell$\to$DBI flows for $p=1,2$.

\section*{1. Principal achievement}
The main achievement is a \textbf{single organising principle} that ties together the matrix
theory decoupling limit, the AdS/CFT near-horizon limit and non-Lorentzian holography.
Under this principle the near-horizon limit is \emph{the same} BPS limit, viewed from the
bulk (Sec.~3.1). Iterating it becomes a mechanical generator of new holographic pairs
(Secs.~3.2--3.3, 4.3--4.4). The individual ingredients are partly known: Polchinski's and
IMSY's treatment of D0-branes, the Maldacena scaling, and the Gomis--Ooguri and
Danielsson--Guijosa--Kruczenski limits. The authors acknowledge this. What is new is that they
handle all of them at once, geometrically and for general $p$, and that the near-horizon metric
explains \emph{why} the boundary of a brane throat is a flat non-Lorentzian M$p$T geometry.
The scattered non-Lorentzian holography proposals of 2024 then appear as entries in one
classification, not as separate ad hoc limits.

The most concrete and checkable result is the family of $p$-brane $T\bar T$ flow equations,
\[
\partial_\lambda\mathcal L=\tfrac{1}{2\lambda^2}\Bigl\{\operatorname{tr}(\mathbb 1-\lambda\mathcal T)
-(p-1)\bigl[\det(\mathbb 1-\lambda\mathcal T)\bigr]^{1/(p-1)}-2\Bigr\}.
\]
For $p=1$ this reduces to ordinary $T\bar T$, and for $p=0$ to the quantum-mechanical flow
of Gross et al. In $d=3$ it gives the strikingly simple membrane flow
$\tfrac14[\operatorname{tr}\mathcal T^2-(\operatorname{tr}\mathcal T)^2+2\lambda\det\mathcal T]$.
The derivation, which reads off the $\omega^{-2}$-dependence of the brane action, is
transparent. I checked numerically that the free-scalar Nambu--Goto Lagrangian
$\lambda^{-1}[1-\sqrt{-\det(\eta+\lambda\,\partial X\!\cdot\!\partial X)}]$ satisfies this
equation, using the paper's stress tensor, for $p=1,\dots,5$ and several numbers of scalars.
The results also settle a puzzle in the literature: why ordinary $T\bar T$ acting on 2d Maxwell
plus scalars does not reach D1 DBI. They also show that the single-scalar ``root'' flows used
elsewhere do not extend to many scalars.

\section*{2. Principal drawback}
The main weakness is that the central conjecture, holography as DLCQ$^n$/DLCQ$^m$, is
\textbf{kinematic and untested}. It is a way of counting and labelling limits on the two
sides of a duality. For none of the new (non-Lorentzian) examples does the paper compute an
observable on both sides: no symmetry algebra, spectrum, correlator, index or thermodynamic
quantity is matched. The existing Lorentzian cases are relabelled rather than re-derived. Several
gaps sharpen this concern, and the authors are open about most of them.
\begin{itemize}
\item The non-Lorentzian bulk backgrounds are \emph{not shown to solve} M$p$T supergravity,
including its torsional constraints. Their isometries and (super)symmetries are not computed
(Outlook).
\item The generic misaligned case ($q\neq p+4$) requires rescaling $\ell$, and hence $N$,
with $\omega$ by hand. The physical meaning of this rescaling is left open.
\item The ``intrinsically non-Lorentzian asymptotic geometry'' lives exactly where the
supergravity description fails. For $p<3$ $\alpha'$ corrections take over, and for $p>3$ the
theory becomes strongly coupled. The analysis is also restricted to the Poincar\'e patch.
\item The ``$T\bar T$ generates the bulk'' statement amounts to rewriting the near-horizon metric
via the emergent dilatation symmetry ($\omega\to\mathbb{\Omega}(r)$). It is classical and does
not make clear what plays the role of the deformation in the field theory. The worldvolume flows
are classical, bosonic and abelian. For $p\ge3$ with gauge fields they are only perturbative.
Nothing is said about whether they share the quantum solvability of 2d $T\bar T$.
\item The proposed DLCQ$^2$/DLCQ$^3$ example (Spin Matrix Theory) depends on the
loosely stated claim that making an R-charge critical ``may be associated'' with a third DLCQ.
\end{itemize}
The paper is therefore a valuable framework and dictionary of limits, but its headline claim
about holography is not yet falsifiable as presented.

\section*{3. Most worthwhile follow-up (feasibility vs.\ impact)}
\textbf{Proposal: turn one DLCQ$^1$/DLCQ$^2$ pair into a tested duality. The best candidate is
the M0T limit of the D4-brane throat (Sec.~3.2 of the paper).}

\emph{Why this example.} On the bulk side the background is explicit: the near-horizon M$p$T background of Sec.~3.2 with
$p=0,\,q=4$, which is dilatation-equivalent to an M0T limit of AdS$_6\times S^4$. It uplifts to the
M5-brane throat AdS$_7\times S^4$ with a null compactification. On the field-theory side the dual
is the DLCQ of the 6d $(2,0)$ theory, i.e.\ quantum mechanics on instanton moduli space.
This is a theory with an independent, long-established definition (Aharony--Berkooz--Kachru--Seiberg--Silverstein)
and with computational tools available: localisation, instanton partition functions and indices.
Both sides are therefore well defined and computable. No other DLCQ$^1$/DLCQ$^2$ example
in the paper has this property.

\emph{Plan.}
\begin{enumerate}[label=(\alph*)]
\item \emph{Bulk consistency (low risk, a few months).} Verify that the background solves the
equations of M0T/non-relativistic M-theory supergravity, including the torsion constraints.
The 11d theory is now available from Bergshoeff et al. Then compute its Killing (super)algebra.
This fills the most glaring gap identified above.
\item \emph{Symmetry matching.} Compare the algebra from (a) with the spacetime symmetry
algebra of the non-relativistic 5d SYM / instanton quantum mechanics found by Lambert et al.
This is the first non-trivial check that the ``bulk = M$p$T limit of string theory''
identification is right.
\item \emph{A protected observable (higher impact).} Match the spectrum of protected
operators, or a suitable index, at fixed null momentum $k$. On one side this means linearised
fluctuations on the non-Lorentzian background. On the other it means the large-$N$ limit of
the known instanton-moduli-space index. A mismatch would also be informative: it would locate
precisely where the DLCQ subtleties (zero modes, non-commuting limits) break the conjecture.
\end{enumerate}

\emph{Why this balances feasibility and impact.} Steps (a) and (b) are classical
calculations with well-defined end-points, and they produce publishable results whatever
happens in (c). Step (c) builds on existing localisation technology rather than
requiring new non-perturbative methods. If the matching works, it would be the first dynamical
evidence for the DLCQ$^n$/DLCQ$^{n+1}$ picture. It would also give a top-down template, with
a holographic dictionary, for the whole family of non-Lorentzian dualities the paper generates.

A cheaper alternative with lower impact is to study the quantum properties of the
$d=3$ flow $\partial_\lambda\mathcal L=\tfrac14[\operatorname{tr}\mathcal T^2-(\operatorname{tr}\mathcal T)^2+2\lambda\det\mathcal T]$,
for example by asking whether it induces a solvable deformation of energy levels on $T^2$. It is
a natural side project, but it is unlikely to bear on the paper's main holographic claims.

\end{document}
